\ncdmodultestmthree{3}{Uddannelse, Fællesskab og Kultur}{全真模考卷三：音乐节义工、好评推荐与社团文化}{音乐节志愿者故事 / 小组报告与邻居代收快递聊天 / 开设自行车咖啡馆 / 三人谈社团生活 / 车行好评推荐与求职面试回信 / 大学学业与志愿工作陈述及消费购物四图双人互动}{\textit{-s} 被动 vs \textit{blive} 被动 · 介词 + 形式宾语 \textit{det} + 从句 · 现在/过去分词定语 · 省略 \textit{hvis} 条件倒装 · 篇章因果与评价转述}

% =====================================================================
% DEL 1: LÆSNING
% =====================================================================
\ncdsubsec{Del 1: Læsning（阅读理解｜50 minutter · 24 point）}

\begin{tipbox}{考场规则与配速建议（Læsning｜15--24 point = Bestået）}
\small
\textit{Læsning opgave 1, 2A, 2B, 3 og 4 udleveres samtidig. Samlet tid: 50 minutter. Kursisten må ikke bruge hjælpemidler.}（阅读五部分合订发放，\textbf{总计时 50 分钟}，全程禁用词典或手机。满分 \textbf{24 分}，\textbf{15--24 分及格（Bestået）}；0--14 分不及格（Ikke bestået）。建议配速：Opgave 1（8 min）$\rightarrow$ Opgave 2A \& 2B（10 min）$\rightarrow$ Opgave 3（14 min）$\rightarrow$ Opgave 4（13 min），最后留 5 分钟复查。）
\end{tipbox}

% ---------------------------------------------------------------------
\subsubsection*{Opgave 1 (6 point)：Frivillig på sommerens musikfestival}

\textbf{Instruktion:}
\begin{itemize}
  \item Læs teksten \textit{Frivillig på sommerens musikfestival}.
  \item Skriv de ord (1--6), der mangler. Du skal kun bruge hvert ord én gang. Der er seks ord, du ikke skal bruge. Se eksemplet (0).
\end{itemize}

\begin{emailbox}{Frivillig på sommerens musikfestival}
Freja elsker livemusik og har altid haft lyst til at opleve en stor dansk musikfestival indefra. I år har hun \textbf{(0)} \underline{\textbf{besluttet}} sig for at arbejde som frivillig på en festival i Roskilde. Det kræver en stor indsats at stå op i mange timer. Freja glæder sig meget til opgaven, \textbf{(1)}\ \rule{1.9cm}{0.4pt} hun ved, at det også bliver hårdt.

\vspace{0.35em}
Den første dag møder Freja ind om morgenen og får sin T-shirt udleveret. Hun skal stå i et stort telt ved indgangen, \textbf{(2)}\ \rule{1.9cm}{0.4pt} alle gæsterne henter deres armbånd. Hun skal også hjælpe folk med at finde deres \textbf{(3)}\ \rule{1.9cm}{0.4pt} på campingområdet, så der ikke opstår lange køer. Det er travlt om eftermiddagen, når tusindvis af gæster ankommer med telte og rygsække. Heldigvis arbejder Freja \textbf{(4)}\ \rule{1.9cm}{0.4pt} med fire andre frivillige, og de hjælper hele tiden hinanden.

\vspace{0.35em}
På den tredje dag begynder det desværre at regne kraftigt, og hendes sko bliver helt våde. Freja overvejer et øjeblik at tage tidligt hjem, \textbf{(5)}\ \rule{1.9cm}{0.4pt} hun fryser og er træt i benene. Men da hendes yndlingsband går på scenen om aftenen, glemmer hun hurtigt regnen og kulden. Til sidst er hun glad for, at hun blev. Vejret om sommeren er ikke altid perfekt, men hun vil \textbf{(6)}\ \rule{1.9cm}{0.4pt} melde sig som frivillig igen næste år.
\end{emailbox}

\begin{grammatikbox}{Ordboks（选词框 · 共 13 词：1 例词 + 6 答案 + 6 多余干扰词）}
\centering
\large\bfseries\color{NCDNavy}
\tikz[baseline=(w.base)]{\node[inner sep=1pt](w){\color{NCDGray}besluttet}; \draw[NCDRed,line width=0.9pt](w.west)--(w.east);} \qquad
selvom \qquad alene \qquad hvor \qquad derfor \qquad plads \qquad heller\\[3pt]
sammen \qquad hvornår \qquad fordi \qquad pause \qquad alligevel \qquad som
\end{grammatikbox}

% ---------------------------------------------------------------------
\subsubsection*{Opgave 2A (NY · 3 point)：Jonas skriver til sin studiegruppeven Maria}

\textbf{Instruktion:}
\begin{itemize}
  \item \textbf{Situation:} Jonas skriver til sin studiegruppeven Maria.
  \item Læs Jonas’ og Marias chat. Der er 5 beskeder fra Jonas.
  \item Find de svar fra Maria, der passer (A--F).
  \item Skriv bogstavet på det svar, der passer. Der er tre svar, du ikke skal bruge. Se eksemplet (0).
\end{itemize}

\small
\renewcommand{\arraystretch}{1.3}
\begin{longtable}{@{} >{\raggedright\arraybackslash}p{7.6cm} @{\hspace{0.4cm}} >{\raggedright\arraybackslash}p{8.8cm} @{}}
\toprule
\textbf{Chatbeskeder mellem Jonas og Maria} & \textbf{Marias mulige svar（Z 为例答，A--F 选 3 项）} \\
\midrule
\endfirsthead
\toprule
\textbf{Chatbeskeder（续表）} & \textbf{Marias mulige svar（续表）} \\
\midrule
\endhead
\bottomrule
\endfoot
\badge{NCDNavy}{Jonas}\ Hej Maria! Har du tid til at arbejde videre på vores oplæg til på torsdag?\newline
\textbf{0.}\ \badge{NCDTeal}{Maria}\ \fbox{\parbox[c][1.1em][c]{1.1em}{\centering\textbf{Z}}} &
\badge{NCDGray}{Z · Eksempel}\ \textcolor{NCDGray}{\textit{Hej Jonas! Ja, det har jeg. Skal vi ikke mødes på biblioteket i morgen formiddag?}} \\
\midrule
\badge{NCDNavy}{Jonas}\ Jo, det er en god idé. Skal vi sige kl. 10 i stueetagen?\newline
\textbf{1.}\ \badge{NCDTeal}{Maria}\ \fbox{\parbox[c][1.1em][c]{1.1em}{\centering\phantom{Z}}} &
\badge{NCDNavy}{A}\ Super, kl. 10 passer mig perfekt. Men jeg har allerede lavet hele oplægget, så vi behøver slet ikke mødes længe.\newline
\badge{NCDNavy}{B}\ Det passer mig rigtig godt med afslutningen. Men måske skal vi booke et grupperum, så der ikke er så meget støj? \\
\midrule
\badge{NCDNavy}{Jonas}\ Fint, så ses vi kl. 10.30. Jeg kan godt tage indledningen og tallene, hvis du vil lave afslutningen?\newline
\textbf{2.}\ \badge{NCDTeal}{Maria}\ \fbox{\parbox[c][1.1em][c]{1.1em}{\centering\phantom{Z}}} &
\badge{NCDNavy}{C}\ Nej, lad den endelig blive hjemme, for vi skal bare læse mine noter igennem på papir.\newline
\badge{NCDNavy}{D}\ Jeg har desværre lægetid om morgenen, så jeg kan først kl. 10.30. Har du tænkt på, hvordan vi skal fordele delene? \\
\midrule
\badge{NCDNavy}{Jonas}\ Jeg skal nok booke et grupperum i dag, så vi kan sidde i fred og ro. Skal jeg også tage min computer med?\newline
\textbf{3.}\ \badge{NCDTeal}{Maria}\ \fbox{\parbox[c][1.1em][c]{1.1em}{\centering\phantom{Z}}} &
\badge{NCDNavy}{E}\ Ja, meget gerne, for min egen er desværre ret langsom for tiden. Så kan vi lettere samle vores slides.\newline
\badge{NCDNavy}{F}\ Jeg vil hellere have, at vi aflyser mødet på biblioteket og bare sender vores tekst til hinanden i aften. \\
\midrule
\badge{NCDNavy}{Jonas}\ Aftale! Så tager jeg den med, og vi ses i morgen kl. 10.30. &
\textcolor{NCDGray}{\small （提示：从 A--F 中选出 3 项填入左侧 1--3 方框，另有 3 项多余。）} \\
\end{longtable}
\normalsize

% ---------------------------------------------------------------------
\subsubsection*{Opgave 2B (NY · 3 point)：Hanne skriver til sin nabo Ali}

\textbf{Instruktion:}
\begin{itemize}
  \item \textbf{Situation:} Hanne skriver til sin nabo Ali.
  \item Læs Hannes og Alis chat. Der er 5 beskeder fra Hanne.
  \item Find de svar fra Ali, der passer (A--F).
  \item Skriv bogstavet på det svar, der passer. Der er tre svar, du ikke skal bruge. Se eksemplet (0).
\end{itemize}

\small
\renewcommand{\arraystretch}{1.3}
\begin{longtable}{@{} >{\raggedright\arraybackslash}p{7.6cm} @{\hspace{0.4cm}} >{\raggedright\arraybackslash}p{8.8cm} @{}}
\toprule
\textbf{Chatbeskeder mellem Hanne og Ali} & \textbf{Alis mulige svar（Z 为例答，A--F 选 3 项）} \\
\midrule
\endfirsthead
\toprule
\textbf{Chatbeskeder（续表）} & \textbf{Alis mulige svar（续表）} \\
\midrule
\endhead
\bottomrule
\endfoot
\badge{NCDNavy}{Hanne}\ Hej Ali! Jeg skal på konference med arbejdet fra onsdag til søndag. Har du mulighed for at hjælpe mig med en lille ting?\newline
\textbf{0.}\ \badge{NCDTeal}{Ali}\ \fbox{\parbox[c][1.1em][c]{1.1em}{\centering\textbf{Z}}} &
\badge{NCDGray}{Z · Eksempel}\ \textcolor{NCDGray}{\textit{Hej Hanne! Ja, selvfølgelig. Hvad kan jeg hjælpe med?}} \\
\midrule
\badge{NCDNavy}{Hanne}\ Tusind tak! Der kommer nemlig en vigtig pakke til pakkeshoppen i Brugsen torsdag, og jeg har også nogle planter på altanen.\newline
\textbf{1.}\ \badge{NCDTeal}{Ali}\ \fbox{\parbox[c][1.1em][c]{1.1em}{\centering\phantom{Z}}} &
\badge{NCDNavy}{A}\ Det gør jeg. Hvornår kommer du egentlig hjem igen på søndag, og hvornår skal jeg aflevere nøglen tilbage?\newline
\badge{NCDNavy}{B}\ Jeg skal desværre selv rejse væk fra torsdag til søndag, så jeg kan ikke passe dine ting alligevel. \\
\midrule
\badge{NCDNavy}{Hanne}\ Du får koden til pakken på sms, og min ekstranøgle ligger i din postkasse i morgen aften. Planterne skal bare have vand fredag.\newline
\textbf{2.}\ \badge{NCDTeal}{Ali}\ \fbox{\parbox[c][1.1em][c]{1.1em}{\centering\phantom{Z}}} &
\badge{NCDNavy}{C}\ Intet problem! Jeg henter gerne pakken og vander dine planter. Hvad skal jeg bruge i pakkeshoppen, og hvordan kommer jeg ind?\newline
\badge{NCDNavy}{D}\ Super, så henter jeg pakken. Du behøver ikke lægge nogen nøgle til mig, for min kone har allerede din ekstranøgle. \\
\midrule
\badge{NCDNavy}{Hanne}\ Den er ret lille og vejer næsten ingenting, så du kan bare lægge den på bordet i min entré.\newline
\textbf{3.}\ \badge{NCDTeal}{Ali}\ \fbox{\parbox[c][1.1em][c]{1.1em}{\centering\phantom{Z}}} &
\badge{NCDNavy}{E}\ Fint, jeg stiller bare pakken ind i min egen stue, indtil du kommer hjem om et par uger.\newline
\badge{NCDNavy}{F}\ Det lyder helt fint. Er pakken stor og tung, og hvor skal jeg stille den, når jeg har hentet den? \\
\midrule
\badge{NCDNavy}{Hanne}\ Jeg lander først søndag aften ved 21-tiden, så jeg henter bare nøglen hos dig mandag aften. Endnu en gang tak! &
\textcolor{NCDGray}{\small （提示：从 A--F 中选出 3 项填入左侧 1--3 方框，另有 3 项多余。）} \\
\end{longtable}
\normalsize

% ---------------------------------------------------------------------
\subsubsection*{Opgave 3 (NY · 5 point)：Søren åbner en lille cykelcafé}

\textbf{Instruktion:}
\begin{itemize}
  \item Læs teksten \textit{Søren åbner en lille cykelcafé}. Der er fem afsnit.
  \item I hvert afsnit mangler der en sætning. Under hvert afsnit står der fire sætninger.
  \item Skriv bogstavet på den sætning, der passer i afsnittet. Der er tre sætninger i hvert afsnit, du ikke skal bruge.
\end{itemize}

\begin{tcolorbox}[enhanced, breakable, colback=NCDBgGray, colframe=NCDNavy!60, boxrule=0.7pt, arc=2pt, left=9pt, right=9pt, top=6pt, bottom=6pt]
\textbf{\color{NCDNavy}1.}\ Søren har arbejdet som cykelmekaniker i Odense i næsten ti år, og han er rigtig dygtig til sit fag. Alligevel har han længe drømt om at prøve noget nyt, hvor han både kan bruge sine hænder og snakke endnu mere med folk i kvarteret. En dag læser han en artikel om cykelcaféer i udlandet, hvor man kan få repareret sin cykel, mens man drikker en god kop kaffe. \rule{2.3cm}{0.4pt} For han har allerede en lille opsparing i banken og kender mange cyklister i byen.\\[4pt]
\textbf{A.} Men Søren synes, at kaffe og cykler slet ikke passer sammen.\\
\textbf{B.} Derfor vælger han at blive ansat på et endnu større værksted.\\
\textbf{C.} Han beslutter sig straks for selv at åbne sådan et sted i Odense.\\
\textbf{D.} Desværre har han aldrig før prøvet at reparere en cykel.
\end{tcolorbox}

\begin{tcolorbox}[enhanced, breakable, colback=NCDBgGray, colframe=NCDNavy!60, boxrule=0.7pt, arc=2pt, left=9pt, right=9pt, top=6pt, bottom=6pt]
\textbf{\color{NCDNavy}2.}\ Det tager dog flere måneder at finde de rigtige lokaler til den nye café. Søren kigger på tre forskellige butikker i centrum, men huslejen er alt for høj til en nystartet forretning. Til sidst finder han et gammelt hjørnelokale på en stille sidegade tæt på en populær cykelsti. \rule{2.3cm}{0.4pt} Lokalet trænger nemlig til ny maling og et bedre gulv, før han kan sætte kaffemaskinen op.\\[4pt]
\textbf{A.} Prisen er heldigvis rimelig, selvom der skal laves en del arbejde indvendigt.\\
\textbf{B.} Derfor åbner han caféen allerede dagen efter uden at ændre noget.\\
\textbf{C.} Men beliggenheden er dårlig, fordi der aldrig kommer cyklister forbi.\\
\textbf{D.} Alligevel er huslejen alt for dyr, så han må droppe drømmen helt.
\end{tcolorbox}

\begin{tcolorbox}[enhanced, breakable, colback=NCDBgGray, colframe=NCDNavy!60, boxrule=0.7pt, arc=2pt, left=9pt, right=9pt, top=6pt, bottom=6pt]
\textbf{\color{NCDNavy}3.}\ Før åbningen får Søren brug for ekstra hjælp til cafédelen, fordi han ikke selv ved så meget om bønner og mælkeskum. Han sætter derfor et opslag op på sociale medier, og allerede næste dag skriver en ung kvinde ved navn Maja, som har arbejdet som barista i fire år. De mødes til en samtale og bliver hurtigt enige om et samarbejde. \rule{2.3cm}{0.4pt} På den måde kan Søren fokusere på dæk, bremser og gear, mens Maja sørger for kaffe og nybagte boller.\\[4pt]
\textbf{A.} Men Maja vil kun arbejde der, hvis hun også må reparere cyklerne.\\
\textbf{B.} Derfor beslutter Søren sig for slet ikke at servere kaffe i butikken.\\
\textbf{C.} Hun har desværre ingen erfaring fra cafébranchen overhovedet.\\
\textbf{D.} Maja skal arbejde deltid og stå for disken, især om morgenen og i weekenden.
\end{tcolorbox}

\begin{tcolorbox}[enhanced, breakable, colback=NCDBgGray, colframe=NCDNavy!60, boxrule=0.7pt, arc=2pt, left=9pt, right=9pt, top=6pt, bottom=6pt]
\textbf{\color{NCDNavy}4.}\ Lørdag den 3. maj slår Søren og Maja endelig dørene op til \textit{Kæde \& Kaffe}. Vejret er solrigt og lunt, og allerede fra klokken ni står de første nysgerrige naboer klar udenfor. Nogle kommer kun for at drikke en latte i solen, mens andre triller ind med en punkteret cykel. \rule{2.3cm}{0.4pt} Det betyder, at stemningen i lokalet er utrolig hyggelig, selvom der er travlt hele dagen.\\[4pt]
\textbf{A.} Folk bliver dog sure over, at de skal vente hele tre minutter på kaffen.\\
\textbf{B.} Mange falder hurtigt i snak med hinanden ved de lange træborde.\\
\textbf{C.} Desværre er gaden helt tom, fordi ingen har set reklamen.\\
\textbf{D.} Derfor lukker de caféen tidligt allerede klokken elleve om formiddagen.
\end{tcolorbox}

\begin{tcolorbox}[enhanced, breakable, colback=NCDBgGray, colframe=NCDNavy!60, boxrule=0.7pt, arc=2pt, left=9pt, right=9pt, top=6pt, bottom=6pt]
\textbf{\color{NCDNavy}5.}\ Her et halvt år efter åbningen går forretningen stadig over al forventning. Hver morgen stopper mange pendlere på vej til arbejde for at få pumpet dækkene og tage en kaffe med på farten. En gang om måneden holder Søren også et lille aftenkursus, hvor gæsterne kan lære at lappe deres egne slanger. \rule{2.3cm}{0.4pt} Søren har derfor aldrig fortrudt, at han sagde sit gamle job op og fulgte sin idé.\\[4pt]
\textbf{A.} Men pendlere vil hellere tage bilen end at cykle i byen.\\
\textbf{B.} Alligevel overvejer han at sælge værkstedet, fordi det er for kedeligt.\\
\textbf{C.} De aftener er næsten altid udsolgte, og flere deltagere kommer igen næste måned.\\
\textbf{D.} Kurserne bliver dog aflyst hver gang, fordi ingen tilmelder sig.
\end{tcolorbox}

% ---------------------------------------------------------------------
\subsubsection*{Opgave 4 (NY · 7 point)：Tre personer fortæller om at være med i en forening}

\textbf{Instruktion:}
\begin{itemize}
  \item Læs de tre tekster (A, B og C), og læs spørgsmålene (1--7).
  \item Find den person, der passer til hvert af de syv spørgsmål.
  \item Sæt kryds (\textbf{X}) ud for den person (A, B eller C), der passer til spørgsmålet. Se eksemplet (0).
\end{itemize}

\small
\renewcommand{\arraystretch}{1.28}
\begin{longtable}{@{} >{\raggedright\arraybackslash}p{10.5cm} @{\hspace{0.25cm}} >{\centering\arraybackslash}p{1.9cm} @{\hspace{0.15cm}} >{\centering\arraybackslash}p{1.9cm} @{\hspace{0.15cm}} >{\centering\arraybackslash}p{1.9cm} @{}}
\toprule
\textbf{Spørgsmål} & \textbf{A. Katrine} & \textbf{B. David} & \textbf{C. Yuna} \\
\midrule
\endfirsthead
\toprule
\textbf{Spørgsmål（续表）} & \textbf{A. Katrine} & \textbf{B. David} & \textbf{C. Yuna} \\
\midrule
\endhead
\bottomrule
\endfoot
0. Hvem har været medlem af sin forening i længst tid? \textcolor{NCDGray}{\small (Eksempel)} & A\ \textcolor{NCDRed}{\textbf{X}} & B \ \underline{\hspace{0.35cm}} & C \ \underline{\hspace{0.35cm}} \\
\midrule
1. Hvem bruger også foreningen til at blive bedre til det danske sprog? & A \ \underline{\hspace{0.35cm}} & B \ \underline{\hspace{0.35cm}} & C \ \underline{\hspace{0.35cm}} \\
\midrule
2. Hvem begyndte i foreningen, fordi personen havde smerter i kroppen? & A \ \underline{\hspace{0.35cm}} & B \ \underline{\hspace{0.35cm}} & C \ \underline{\hspace{0.35cm}} \\
\midrule
3. Hvem var nødt til at lære bestemte regler og teknikker, før personen måtte deltage fuldt ud? & A \ \underline{\hspace{0.35cm}} & B \ \underline{\hspace{0.35cm}} & C \ \underline{\hspace{0.35cm}} \\
\midrule
4. Hvem oplever stor aldersforskel blandt medlemmerne i sin forening? & A \ \underline{\hspace{0.35cm}} & B \ \underline{\hspace{0.35cm}} & C \ \underline{\hspace{0.35cm}} \\
\midrule
5. Hvem optræder af og til for et publikum? & A \ \underline{\hspace{0.35cm}} & B \ \underline{\hspace{0.35cm}} & C \ \underline{\hspace{0.35cm}} \\
\midrule
6. Hvem mødes med sin forening tidligt om morgenen på hverdage? & A \ \underline{\hspace{0.35cm}} & B \ \underline{\hspace{0.35cm}} & C \ \underline{\hspace{0.35cm}} \\
\midrule
7. Hvem fortæller, at det er ret billigt at være med, fordi folk selv hjælper til? & A \ \underline{\hspace{0.35cm}} & B \ \underline{\hspace{0.35cm}} & C \ \underline{\hspace{0.35cm}} \\
\end{longtable}
\normalsize

\begin{center}
\textbf{\large\color{NCDNavy} Tre personer fortæller om at være med i en forening}
\end{center}

\begin{emailbox}{A. Katrine}
Jeg meldte mig ind i vores lokale naturforening for tre år siden, kort tid efter at jeg var flyttet til Silkeborg. Jeg kendte næsten ingen i byen, og i stedet for at sidde alene i min lejlighed ville jeg gerne ud i skoven sammen med andre. Hver anden søndag mødes vi en gruppe på femten til tyve personer og går en tur på otte til ti kilometer. Nogle gange har vi en biolog med, som fortæller om planter og fugle undervejs, og bagefter laver vi varm suppe eller kaffe over bål. Medlemmerne er mellem tyve og halvfjerds år, så man snakker med mange forskellige mennesker. Det koster kun tre hundrede kroner om året, fordi alle skiftes til at planlægge turene frivilligt. Det har været den bedste måde for mig at lære både naturen og området at kende på.
\end{emailbox}

\begin{emailbox}{B. David}
For halvandet år siden begyndte jeg i en roklub ved fjorden, fordi jeg havde fået ondt i ryggen af at sidde foran en skærm hele dagen. I starten var jeg nødt til at tage et obligatorisk begynderkursus over fire aftener, hvor vi lærte teknik og sikkerhed på vandet, før vi måtte ro ud alene. Det var faktisk ret hårdt for armene i begyndelsen! Nu ror jeg sammen med et fast hold på fire mand hver tirsdag og torsdag morgen allerede klokken halv syv, før jeg møder på arbejde. Vandet er helt stille på det tidspunkt, og det giver mig masser af energi til resten af dagen. Kontingentet er lidt højt, omkring to tusinde kroner om året, men til gengæld er alle både og redningsveste med i prisen. Om sommeren tager klubben også på en weekendtur til en anden ø.
\end{emailbox}

\begin{emailbox}{C. Yuna}
Da jeg kom til Aarhus fra Japan for to år siden, savnede jeg at lave noget kreativt efter arbejde og samtidig øve mit dansk i en afslappet stemning. En kollega tog mig med til en øveaften i et lokalt kor, og jeg blev bidt af det med det samme. Man behøver ikke at kunne læse noder eller gå til optagelsesprøve for at være med, så alle kan starte. Vi øver hver onsdag aften i et kulturhus, hvor vi synger både danske sange og engelske popnumre. På mine første øveaftener var det lidt svært at følge med i de hurtige danske tekster, men de andre medlemmer hjalp mig tålmodigt med udtalen. To gange om året giver vi små koncerter på plejehjem eller i en kirke. For mig giver koret et stærkt fællesskab, og flere af sangerne er i dag mine nære venner.
\end{emailbox}

% =====================================================================
% DEL 2: SKRIVNING
% =====================================================================
\ncdsubsec{Del 2: Skrivning（书面表达｜Opgave 1 \& Opgave 2）}

\subsubsection*{Opgave 1（半正式好评推荐信 · Anbefaling）}

\textbf{Situation:} Du har fået repareret din cykel hos den lokale cykelhandler \textit{”Byens Cykler“}, og du er meget tilfreds med servicen. Du vil skrive en anbefaling på internettet.\\
\textbf{Opgave:} Skriv en anbefaling af \textit{”Byens Cykler“}. Du skal fortælle:
\begin{itemize}
  \item hvor værkstedet ligger
  \item hvad du fik lavet
  \item hvorfor du synes, servicen var god
  \item hvem du vil anbefale værkstedet til
\end{itemize}

\begin{emailbox}{Anbefaling af ”Byens Cykler“}
\vspace{0.25em}
\rule{\linewidth}{0.4pt}\\[0.52em]
\rule{\linewidth}{0.4pt}\\[0.52em]
\rule{\linewidth}{0.4pt}\\[0.52em]
\rule{\linewidth}{0.4pt}\\[0.52em]
\rule{\linewidth}{0.4pt}\\[0.52em]
\rule{\linewidth}{0.4pt}\\[0.52em]
\rule{\linewidth}{0.4pt}\\[0.52em]
\rule{\linewidth}{0.4pt}
\end{emailbox}

\subsubsection*{Opgave 2（私人回信 · Minimum 90 ord · Tid: 30 minutter）}

\textbf{Situation:} Du har fået e-mailen her nedenfor fra din ven, Henrik.\\
\textbf{Opgave:} Læs e-mailen, og skriv et svar til Henrik. Du skal svare på alle Henriks spørgsmål, og du skal indlede og afslutte e-mailen på en passende måde.\\
\textbf{Længde:} Minimum 90 ord. \qquad \textbf{Tid:} 30 minutter.

\begin{emailbox}{Modtaget e-mail fra Henrik}
\hfill Aarhus d. 14. oktober 2025

Kære ....

Hvordan går det? Jeg hørte fra din søster, at du har været til en spændende jobsamtale for nylig! Tillykke med at du kom til samtale! Jeg er jo selv ved at søge arbejde, så jeg vil rigtig gerne høre om det.

Hvordan gik det til samtalen? Hvad handler jobbet helt præcist om? Hvordan var selve samtalen og personerne, du talte med? Og hvordan forberedte du dig på forhånd på dansk?

Jeg glæder mig meget til at høre fra dig!

Mange hilsner\\
Henrik
\end{emailbox}

\begin{emailbox}{Dit svar｜Til: Henrik \qquad Emne: Jobsamtalen}
\vspace{0.25em}
\rule{\linewidth}{0.4pt}\\[0.52em]
\rule{\linewidth}{0.4pt}\\[0.52em]
\rule{\linewidth}{0.4pt}\\[0.52em]
\rule{\linewidth}{0.4pt}\\[0.52em]
\rule{\linewidth}{0.4pt}\\[0.52em]
\rule{\linewidth}{0.4pt}\\[0.52em]
\rule{\linewidth}{0.4pt}\\[0.52em]
\rule{\linewidth}{0.4pt}
\end{emailbox}

% =====================================================================
% DEL 3: MUNDTLIG KOMMUNIKATION
% =====================================================================
\ncdsubsec{Del 3: Mundtlig kommunikation（口语交际全真题卡｜Opgave 1 \& 2）}

\subsubsection*{Mundtlig Opgave 1 (Emneopgave med Mindmap｜Del 1: 1--2 min. + Del 2: 3--4 min.)}

\begin{tipbox}{考场规则提醒（Opgave 1 · 官方红线与陈述要求）}
\small
考生需自带\textbf{两张}提前准备好的个人主题思维导图（至少 1 张须属于工作 / 实习 / 教育 / 丹麦语学习领域）。导图格内\textbf{只能填写关键词或短语（stikord），严禁写完整句子}！官方明确禁止使用复杂度过低的题目（如 \textit{Min bolig}、\textit{Min familie}、\textit{Min hverdag}）。抽签决定 1 张后，先进行 \textbf{1--2 分钟}连贯陈述（Del 1），再接受考官 \textbf{3--4 分钟}追问并简要说明理由（Del 2）。
\end{tipbox}

\begin{grammatikbox}{Mindmap 3A（已填关键词卡）：Min skole / uddannelse}
\small
\renewcommand{\arraystretch}{1.35}
\begin{tabularx}{\linewidth}{@{} >{\leavevmode\bfseries\color{NCDNavy}\raggedright\arraybackslash}p{3.4cm} @{\hspace{0.2cm}} >{\raggedright\arraybackslash}X @{\hspace{0.4cm}} >{\leavevmode\bfseries\color{NCDNavy}\raggedright\arraybackslash}p{3.4cm} @{\hspace{0.2cm}} >{\raggedright\arraybackslash}X @{}}
\toprule
\multicolumn{4}{c}{\large\bfseries\color{NCDRed} Emne: Min skole / uddannelse \quad\small\color{NCDGray}（仅含 stikord 关键词，不含整句）} \\
\midrule
1. sted / tid & København · erhvervsakademi · 2. år · man.--fre. kl. 9--15 &
2. skolen / uddannelsen & offentlig · ca. 200 studerende på årgangen · multimediedesign \\
\midrule
3. transport & cykel 20 min. langs kanalen · metro når det regner &
4. undervisningen & grafisk design · kodning · gruppeprojekter · bedst: praksis / sværest: danske rapporter \\
\midrule
5. lærere / medstuderende & uformelle lærere · danskere og internationale · fast studiegruppe &
6. mening om uddannelsen & rigtig glad · jobrelevant · godt fællesskab · travlt før eksamen \\
\bottomrule
\end{tabularx}
\end{grammatikbox}

\textbf{Del 2 考官追问模拟（Min skole / uddannelse · 4 题）：}
\begin{enumerate}
  \item \textit{Du siger, at I laver mange gruppeprojekter --- kan du bedst lide at arbejde alene eller i en gruppe? Hvorfor?}
  \item \textit{Du nævner, at skriftlige rapporter på dansk er det sværeste --- hvad gør du for at blive bedre til at skrive på dansk?}
  \item \textit{Hvordan er forholdet mellem lærere og studerende på din uddannelse i forhold til dit hjemland?}
  \item \textit{Hvilken slags job drømmer du om at få, når du er færdig med uddannelsen?}
\end{enumerate}

\begin{grammatikbox}{Mindmap 3B（已填关键词卡）：Mit frivillige arbejde}
\small
\renewcommand{\arraystretch}{1.35}
\begin{tabularx}{\linewidth}{@{} >{\leavevmode\bfseries\color{NCDNavy}\raggedright\arraybackslash}p{3.4cm} @{\hspace{0.2cm}} >{\raggedright\arraybackslash}X @{\hspace{0.4cm}} >{\leavevmode\bfseries\color{NCDNavy}\raggedright\arraybackslash}p{3.4cm} @{\hspace{0.2cm}} >{\raggedright\arraybackslash}X @{}}
\toprule
\multicolumn{4}{c}{\large\bfseries\color{NCDRed} Emne: Mit frivillige arbejde \quad\small\color{NCDGray}（仅含 stikord 关键词，不含整句）} \\
\midrule
1. sted / organisation & kulturhus i midtbyen · folkekøkken ”Naboens Bord“ &
2. tidsforbrug & hver anden torsdag aften · kl. 16.30--20.30 · ca. 8 timer/md. \\
\midrule
3. arbejdsopgaver & skære grøntsager · dække langborde · tage imod gæster · rydde op &
4. mennesker / hold & 10 faste frivillige · alle aldre · naboer og studerende · god holdånd \\
\midrule
5. udbytte / dansk & øver hverdagssnak · nye danske venner · lærer om madkultur &
6. mening om frivilligt arbejde & meningsfuldt · giver energi · lidt hårdt for benene · stolt af indsatsen \\
\bottomrule
\end{tabularx}
\end{grammatikbox}

\textbf{Del 2 考官追问模拟（Mit frivillige arbejde · 4 题）：}
\begin{enumerate}
  \item \textit{Hvordan fandt du frem til lige præcis dette folkekøkken, og hvorfor valgte du at blive frivillig?}
  \item \textit{Hvad kan du bedst lide af dine opgaver i køkkenet eller ved bordene, og er der noget, der er mindre sjovt?}
  \item \textit{Hvad har det frivillige arbejde betydet for dit danske sprog og dit netværk i byen?}
  \item \textit{Synes du, at det er let eller svært at passe frivilligt arbejde ved siden af hverdag og arbejde?}
\end{enumerate}

% ---------------------------------------------------------------------
\subsubsection*{Mundtlig Opgave 2 (Paropgave Del 1: ca. 4 min. + Individuel Del 2: ca. 3 min. pr. kursist)}

\begin{grammatikbox}{Kursistark（Opgaveark · EMNE: Indkøb og forbrug）}
\small
\textbf{\color{NCDRed}Opgave:} \textit{Testeren introducerer dig til opgaven.} \qquad \textbf{\color{NCDNavy}Emne:} \textit{Indkøb og forbrug.}\\[4pt]
\renewcommand{\arraystretch}{1.35}
\begin{tabularx}{\linewidth}{@{} >{\raggedright\arraybackslash}X @{\hspace{0.4cm}} >{\raggedright\arraybackslash}X @{}}
\toprule
\badge{NCDNavy}{[图片 1: 本地超市常规采购]}\newline
\textit{En kunde med en indkøbskurv blandt hylder og gule tilbudsskilte i et almindeligt dansk supermarked.} &
\badge{NCDTeal}{[图片 2: 网上购物与送货]}\newline
\textit{En person bestiller dagligvarer og tøj på sin bærbare computer, mens et par pakker allerede står ved hoveddøren.} \\
\midrule
\badge{NCDGold}{[图片 3: 跳蚤市场与二手店]}\newline
\textit{Besøgende kigger på brugte møbler, bøger, lamper og tøj i boderne på et udendørs loppemarked eller i en genbrugsbutik.} &
\badge{NCDRed}{[图片 4: 农场直购/本地集市]}\newline
\textit{Et lokalt torvemarked eller en gårdbutik med friske økologiske grøntsager, æg og brød direkte fra producenten.} \\
\bottomrule
\end{tabularx}
\end{grammatikbox}

\begin{emailbox}{Testerark（考官引导语与 Del 1 / Del 2 考核指令）}
\small
\textbf{\color{NCDNavy}Del 1 (Paropgave -- ca. 4 min.)}\\
Opgavearket lægges foran kursisterne, der præsenteres for emnet på følgende måde:\\[2pt]
\textit{”Emnet er ’Indkøb og forbrug’. Her er nogle billeder, der viser forskellige måder at købe ind på. I skal tale med hinanden om, hvad I synes om de forskellige måder at handle på, og hvorfor. I skal stille spørgsmål til hinanden og svare på hinandens spørgsmål.“}\\[4pt]
Hvis kursisterne ikke udtrykker deres præferencer i forhold til alle fire billeder, må tester sørge for, at dette sker ved at spørge:\\[2pt]
\textit{”Hvad synes I om at handle i et almindeligt supermarked / købe ind på nettet / handle på loppemarked eller i genbrugsbutik / købe varer i en gårdbutik eller på et lokalt marked? Hvorfor? / Hvorfor ikke?“}\\[6pt]
\textbf{\color{NCDNavy}Del 2 (Individuel opgave -- ca. 3 min. pr. kursist)}\\
Tester stiller hver kursist nedenstående spørgsmål:
\begin{enumerate}
  \item \textit{Hvordan kan du bedst lide at købe ind --- og hvorfor?}
  \item \textit{Hvad er vigtigst for dig, når du handler i hverdagen (fx pris, kvalitet, tid, klima eller økologi)?}
  \item \textit{Kan du fortælle lidt om et godt eller dårligt køb, du har lavet for nylig? (hvad, hvor, hvornår, hvorfor godt/dårligt)}
\end{enumerate}
\end{emailbox}

% =====================================================================
% DEL 4: RETTEARK, DETALJERET ANALYSE & HØJSCORE-MODELLER
% =====================================================================
\ncdsubsec{Del 4: Retteark, Detaljeret Analyse \& Højscore-Modeller}

\subsubsection*{4.1 官方 Retteark 评分卡（Læsning 24 point \& Skrivning / Mundtlig 评估基准）}

\small
\renewcommand{\arraystretch}{1.28}
\begin{longtable}{@{} >{\leavevmode\bfseries\color{NCDNavy}\raggedright\arraybackslash}p{3.3cm} @{\hspace{0.25cm}} >{\raggedright\arraybackslash}p{9.45cm} @{\hspace{0.25cm}} >{\centering\arraybackslash}p{1.8cm} @{\hspace{0.25cm}} >{\centering\arraybackslash}p{1.7cm} @{}}
\toprule
\textbf{模块 / 题号} & \textbf{Rettenøgle（标准答案）} & \textbf{分值} & \textbf{自测得分} \\
\midrule
\endfirsthead
\toprule
\textbf{模块 / 题号（续表）} & \textbf{Rettenøgle（标准答案）} & \textbf{分值} & \textbf{自测得分} \\
\midrule
\endhead
\bottomrule
\endfoot
\textbf{Læsning Opgave 1}\newline（6 题） &
1. \textbf{selvom} \quad 2. \textbf{hvor} \quad 3. \textbf{plads} \quad 4. \textbf{sammen} \quad 5. \textbf{fordi} \quad 6. \textbf{alligevel}\newline
\textcolor{NCDGray}{\small （例词 0：besluttet｜多余干扰词：alene, derfor, heller, hvornår, pause, som）} & 6 point & \rule{0.75cm}{0.4pt} / 6 \\
\midrule
\textbf{Læsning Opgave 2A}\newline（3 题） &
1. \textbf{D} \quad 2. \textbf{B} \quad 3. \textbf{E}\newline
\textcolor{NCDGray}{\small （例答 0：Z｜多余干扰项：A, C, F）} & 3 point & \rule{0.75cm}{0.4pt} / 3 \\
\midrule
\textbf{Læsning Opgave 2B}\newline（3 题） &
1. \textbf{C} \quad 2. \textbf{F} \quad 3. \textbf{A}\newline
\textcolor{NCDGray}{\small （例答 0：Z｜多余干扰项：B, D, E）} & 3 point & \rule{0.75cm}{0.4pt} / 3 \\
\midrule
\textbf{Læsning Opgave 3}\newline（5 题） &
1. \textbf{C} \quad 2. \textbf{A} \quad 3. \textbf{D} \quad 4. \textbf{B} \quad 5. \textbf{C} & 5 point & \rule{0.75cm}{0.4pt} / 5 \\
\midrule
\textbf{Læsning Opgave 4}\newline（7 题） &
1. \textbf{C} \quad 2. \textbf{B} \quad 3. \textbf{B} \quad 4. \textbf{A} \quad 5. \textbf{C} \quad 6. \textbf{B} \quad 7. \textbf{A}\newline
\textcolor{NCDGray}{\small （例题 0 为 A）} & 7 point & \rule{0.75cm}{0.4pt} / 7 \\
\midrule
\textbf{Læsning 总计} & \textbf{Bestået：15--24 point（及格）} \quad|\quad \textbf{Ikke bestået：0--14 point} & \textbf{24 point} & \rule{0.75cm}{0.4pt} / 24 \\
\end{longtable}
\normalsize

\begin{tipbox}{Skrivning \& Mundtlig Opgave 1--2 整体评估自检准绳（Bestået 底线）}
\small
\begin{itemize}
  \item \textbf{Skrivning Opgave 1（Anbefaling）}：\ding{172} 必须完整涵盖指定 4 个要点（地点、维修项目、为何满意、推荐给谁）；\ding{173} 语气礼貌得体，灵活运用 \textit{som / fordi / selvom} 与过去时叙事 + 结尾推荐句型（如 \textit{Jeg vil varmt anbefale...}）。
  \item \textbf{Skrivning Opgave 2（Privat e-mail）}：\ding{172} 字数硬性要求 \textbf{minimum 90 ord}（建议写到 105--135 词）；\ding{173} 回答 Henrik 的全部 4 个问题；\ding{174} 具备地道开头（\textit{Kære Henrik}）与落款（\textit{Kærlig/Mange hilsner}）；\ding{175} 保持正确的 B1 从句语序与时态呼应。
  \item \textbf{Mundtlig Opgave 1 \& 2}：Opgave 1 仅凭 \textit{stikord} 连贯陈述 1--2 分钟，并在 Del 2 给出带简短理由（\textit{fordi / derfor}）的回应；Opgave 2 Del 1 两人必须交互问答，对\textbf{全部 4 张图}都明确表达偏好与原因，Del 2 能围绕个人消费习惯展开 3 分钟问答。
\end{itemize}
\end{tipbox}

% ---------------------------------------------------------------------
\subsubsection*{4.2 阅读 Opgave 1--4 逐题深度解析、全文精译与核心词表}

\paragraph*{一、Opgave 1（音乐节志愿者）完形语法逻辑解析与全文精译}
\begin{itemize}
  \item \textbf{(1) selvom}：前后形成“虽然知道辛苦，但依然非常期待”的让步关系（\textit{glæder sig meget..., selvom hun ved...}）。干扰词 \textit{fordi} 逻辑反向（不会“因为工作辛苦而期待”）；\textit{derfor} 属于副词不能直接作从句连词，且后接主句需倒装。
  \item \textbf{(2) hvor}：修饰地点名词 \textit{et stort telt ved indgangen}，引导定语从句说明“在那个帐篷里所有客人都来领取手环”。干扰词 \textit{som} 不能作地点状语（从句中 \textit{alle gæsterne} 和 \textit{deres armbånd} 已完整）；\textit{hvornår} 为疑问副词，不可作定语从句关系词。
  \item \textbf{(3) plads}：放在物主代词 \textit{deres} 之后充当单数共性名词（\textit{finde deres plads på campingområdet}，“找到他们在露营区的营位”）。唯一另一个名词干扰词 \textit{pause}（“找到他们的休息时间”）语义不通。
  \item \textbf{(4) sammen}：构成固定搭配 \textit{arbejde sammen med...}（“与……一起工作”）。干扰词 \textit{alene}（“独自”）与后面的 \textit{med fire andre} 及 \textit{hjælper hinanden} 直接矛盾。
  \item \textbf{(5) fordi}：解释前文“考虑提前回家”的原因（“因为她冷得发抖而且双腿很累”），从句语序 \textit{hun fryser} 正确。干扰词 \textit{derfor} 因果倒置且需触发 V2 倒装（*\textit{derfor fryser hun}）。
  \item \textbf{(6) alligevel}：放在情态动词 \textit{vil} 与实义动词 \textit{melde} 之间，呼应前半句转折连词 \textit{men}（“夏天天气未必完美，但她\textbf{还是}想明年再来当志愿者”）。干扰词 \textit{heller} 只能用于否定句（需配合 \textit{ikke}）。
\end{itemize}

\begin{transbox}
\small
\textbf{\color{NCDTeal}[Opgave 1 全文中文精译]}\\
弗蕾娅热爱现场音乐，一直很想亲身体验一次丹麦大型音乐节的幕后生活。今年她\textbf{（0）决定}去罗斯基勒的一场音乐节做志愿者。这需要付出很大精力、长时间站立。不过弗蕾娅非常期待这项任务，\textbf{（1）尽管}她知道这也会很辛苦。第一天，弗蕾娅一早到岗并领到了自己的 T 恤。她要守在入口处的一个大帐篷里，\textbf{（2）在那里}所有观众会领取他们的手环。她还得帮大家在露营区找到各自的\textbf{（3）营位}，以免排起长队。下午当成千上万名观众背着帐篷和双肩包同时涌入时，工作非常忙碌。幸好弗蕾娅是和另外四名志愿者\textbf{（4）一起}工作，他们一直互相照应。到了第三天，不幸下起了大雨，她的鞋子全湿透了。弗蕾娅有那么一瞬间考虑提前回家，\textbf{（5）因为}她既冷又觉得腿酸。但当她最喜欢的乐队在晚上登台时，她很快便忘记了雨水和寒冷。到最后，她很庆幸自己留了下来。夏天的天气并不总是完美，但她明年\textbf{（6）依然}想再报名当志愿者。
\end{transbox}

\paragraph*{二、Opgave 2A \& 2B（双对话气泡补全）跨回合衔接解析与全文精译}
\begin{itemize}
  \item \textbf{Opgave 2A（Jonas \& Maria）}：
    \begin{itemize}
      \item \textbf{1. D}：Jonas 提议十点在一楼见面；下一句 Jonas 回应 \textit{Fint, så ses vi kl. 10.30. Jeg kan godt tage indledningen... hvis du vil lave afslutningen?}，说明 Maria 必先将时间改到 10:30 并询问“如何分工”（\textit{fordele delene}）。\textbf{排除 A}（A 说 10 点合适且自己写完了，与 10:30 及分工矛盾）。
      \item \textbf{2. B}：承接 Jonas 询问“你做结尾部分可以吗？”，Maria 在 B 中先接受分工（\textit{Det passer mig rigtig godt med afslutningen}），再提议预订小组讨论室以免太吵；Jonas 紧接着在下一条回复 \textit{Jeg skal nok booke et grupperum i dag}，精准扣合。
      \item \textbf{3. E}：回应 Jonas 问“我也要带电脑吗？”，Maria 在 E 中表示“非常欢迎，因为我自己的电脑最近很慢”，且便于合并幻灯片；最后 Jonas 确认 \textit{Så tager jeg den med}。\textbf{排除 C}（拒绝带电脑，与结尾句直接冲突）与 \textbf{F}（提议取消图书馆会面，与订讨论室及明天见矛盾）。
    \end{itemize}
  \item \textbf{Opgave 2B（Hanne \& Ali）}：
    \begin{itemize}
      \item \textbf{1. C}：Hanne 说明有两件事（周四去 Brugsen 包裹点取件 + 给阳台植物浇水）；下一条 Hanne 回复了“短信发取件码、备用钥匙放信箱、周五浇水”，可见 Ali 在 C 中不仅一口答应，还问了“包裹点需要什么”和“怎么进你家”。\textbf{排除 B}（自己要出去旅行）与 \textbf{D}（自称已经有钥匙）。
      \item \textbf{2. F}：承接 Hanne 交代取件码与钥匙后，Ali 在 F 中询问“包裹又大又重吗？取完放哪儿？”，紧接着 Hanne 回复 \textit{Den er ret lille og vejer næsten ingenting, så du kan bare lægge den på bordet i min entré}，问与答严丝合缝。\textbf{排除 E}（说放自己客厅等几周，与放进 Hanne 玄关且周日回来矛盾）。
      \item \textbf{3. A}：Ali 先答应照办（\textit{Det gør jeg}），再问“你周日几点回来、我什么时候还钥匙？”，正好引出 Hanne 最后一条交代周日晚九点落地及周一取钥匙的回复。
    \end{itemize}
\end{itemize}

\begin{transbox}
\small
\textbf{\color{NCDTeal}[Opgave 2A \& 2B 对话中文精译]}\\
\textbf{2A（乔纳斯与玛丽亚）}：乔纳斯：嗨玛丽亚！你有空继续准备周四的报告吗？$\rightarrow$ 玛丽亚（Z 例）：嗨乔纳斯！有的。我们明天上午在图书馆见怎么样？$\rightarrow$ 乔纳斯：好主意。约十点在一楼好吗？$\rightarrow$ 玛丽亚（1. D）：抱歉我早上有医生预约，十点半才能到。你想过怎么分配各部分了吗？$\rightarrow$ 乔纳斯：行，那十点半见。我可以负责开头和数据，你来做结尾好吗？$\rightarrow$ 玛丽亚（2. B）：我做结尾很合适。不过也许我们该订个小组讨论室，这样不会那么吵？$\rightarrow$ 乔纳斯：我今天去订讨论室，好让我们安静坐着。我还要带电脑吗？$\rightarrow$ 玛丽亚（3. E）：太好了，我自己的电脑最近特别慢。这样我们也更好合并幻灯片。$\rightarrow$ 乔纳斯：说定了！我把电脑带上，明天十点半见。\\[3pt]
\textbf{2B（汉娜与阿里）}：汉娜：嗨阿里！我周三到周日要出差开会，能帮我个小忙吗？$\rightarrow$ 阿里（Z 例）：嗨汉娜！当然可以，有什么能帮忙的？$\rightarrow$ 汉娜：太感谢了！周四有个重要包裹会到 Brugsen 的取件点，我阳台上还有一些花。$\rightarrow$ 阿里（1. C）：没问题！我乐意帮忙取件和浇花。去取件需要什么信息？我怎么进你公寓？$\rightarrow$ 汉娜：取件码会发你短信，我的备用钥匙明晚放你信箱。植物周五浇水就行。$\rightarrow$ 阿里（2. F）：听起来很容易。包裹又大又重吗？取回来放哪儿？$\rightarrow$ 汉娜：它挺小的而且几乎不沉，放我玄关桌子上就行。$\rightarrow$ 阿里（3. A）：好的。你周日到底几点回来？我什么时候还钥匙？$\rightarrow$ 汉娜：我周日晚上九点左右落地，周一晚上我去你那儿拿钥匙就行。再次感谢！
\end{transbox}

\paragraph*{三、Opgave 3（开设自行车咖啡馆）段落逻辑补句解析与全文精译}
\begin{itemize}
  \item \textbf{第 1 段 $\rightarrow$ C}：前句写到 Søren 读到国外的自行车咖啡馆报道；后句以 \textit{For} 引出原因“因为他银行里已经有一笔小额存款，而且认识城里很多骑车人”。C 项“他立刻决定自己在欧登塞也开这样一家店”完美构成因果链条（\textbf{排除 A/B/D}：均否定开店或否定十年修车经验）。
  \item \textbf{第 2 段 $\rightarrow$ A}：前句说明他终于在安静小街且靠近热门自行车道旁找到一间旧街角店面；后句用 \textit{nemlig} 解释“店铺确实需要重新粉刷并修整地板才能摆上咖啡机”。A 项“幸好租金合理，尽管内部还得做不少装修”承上启下，\textit{arbejde indvendigt} 直接对应后句的刷墙铺地。
  \item \textbf{第 3 段 $\rightarrow$ D}：前句写咖啡师 Maja 来面试并达成合作意向；后句以 \textit{På den måde} 总结“这样一来，Søren 可以专注轮胎、刹车和变速，而 Maja 负责咖啡和现烤小面包”。D 项“Maja 将做兼职并负责吧台，尤其是早间和周末”准确补出具体的分工方案。
  \item \textbf{第 4 段 $\rightarrow$ B}：前句描述开业当天阳光明媚，有的客人只喝拿铁，有的推着爆胎的自行车进来；后句总结 \textit{Det betyder, at stemningen i lokalet er utrolig hyggelig}（“这意味着店里气氛格外温馨舒适”）。B 项“许多人在长木桌旁很快就聊了起来”正好解释了为何气氛温馨。
  \item \textbf{第 5 段 $\rightarrow$ C}：空前一句提到 Søren 每月还举办一次教客人自己补胎的晚间小课程；后句总结“因此 Søren 从未后悔辞掉旧工作”。C 项中的回指短语 \textit{De aftener}（“那些夜晚”）精准指代前文的 \textit{et lille aftenkursus}，并交代场场爆满的盛况。
\end{itemize}

\begin{transbox}
\small
\textbf{\color{NCDTeal}[Opgave 3 全文中文精译]}\\
\textbf{1.} 索伦在欧登塞当了近十年自行车修理工，手艺非常精湛。尽管如此，他一直梦想尝试点新事物，既能动手干活，又能和街坊邻里多聊聊天。一天，他读到一篇介绍国外“自行车咖啡馆”的文章——在那里人们可以一边喝杯好咖啡，一边修好自行车。\textbf{[C. 他立刻决定自己在欧登塞也开这样一家店。]} 因为他银行里已经有一笔小额存款，而且认识城里许多骑车人。\textbf{2.} 不过找到合适的咖啡馆铺面花了好几个月。索伦看了市中心的三家店铺，但租金对新创小店来说太高了。最后，他在一条靠近热门自行车道的安静支路上找到一间老街角铺面。\textbf{[A. 幸好价格合理，尽管内部还得做不少工程。]} 因为店里在安装咖啡机前确实需要重新粉刷和铺设更好的地板。\textbf{3.} 开业前，索伦在咖啡区需要额外帮手，因为他自己对咖啡豆和打奶泡了解不多。于是他在社交媒体发了招聘帖，第二天就有一位做了四年咖啡师的年轻姑娘玛雅联系他。两人面谈后很快达成合作。\textbf{[D. 玛雅将兼职负责吧台，尤其是早上和周末。]} 这样一来，索伦可以专心处理轮胎、刹车和变速器，而玛雅负责咖啡和现烤面包。\textbf{4.} 五月三日周六，索伦和玛雅终于推开了“链条与咖啡”的大门。天气晴朗温暖，早上九点不到第一批好奇的邻居就等在门外。有人只是来晒太阳喝拿铁，也有人推着扎破胎的车进来。\textbf{[B. 许多人在长木桌旁很快聊了起来。]} 这使得店里的氛围虽然全天忙碌却无比温馨。\textbf{5.} 开业半年后，生意依然远超预期。每天早上都有不少上班族在通勤路上停下来打气并顺手买杯外带咖啡。索伦每月还办一次教客人自己补内胎的晚间小课堂。\textbf{[C. 那些夜晚几乎场场报满，还有好几位学员下个月会再来。]} 因此索伦从未后悔过辞去旧工作、追随自己的想法。
\end{transbox}

\paragraph*{四、Opgave 4（三人谈社团生活）同义定位解析与三人自述精译}
\begin{itemize}
  \item \textbf{(0) A [Eksempel]}：Katrine 三年前入会（\textit{for tre år siden}），David 一年半前（\textit{halvandet år siden}），Yuna 两年前（\textit{to år siden}），故 Katrine 入会时间最长。
  \item \textbf{1. C (Yuna)}：明确提到加入合唱团是为了同时在轻松氛围里练习丹麦语（\textit{samtidig øve mit dansk}），且其他成员耐心帮她纠正发音。
  \item \textbf{2. B (David)}：因为整天坐在电脑屏幕前导致背痛（\textit{havde fået ondt i ryggen}），才加入赛艇俱乐部。对应题干 \textit{havde smerter i kroppen}。
  \item \textbf{3. B (David)}：必须先完成四晚的必修入门课程学习划桨技巧与水上安全，才能单独划出去（\textit{obligatorisk begynderkursus... lærte teknik og sikkerhed... før vi måtte ro ud alene}）。Yuna 明确相反（无需识谱或试音）。
  \item \textbf{4. A (Katrine)}：成员年龄跨度从二十岁到七十岁（\textit{mellem tyve og halvfjerds år}），对应题干 \textit{stor aldersforskel}。
  \item \textbf{5. C (Yuna)}：每年两次在养老院或教堂举办小型音乐会（\textit{giver vi små koncerter}），对应题干 \textit{optræder ... for et publikum}。
  \item \textbf{6. B (David)}：每周二、周四早上六点半上班前一起划船（\textit{tirsdag og torsdag morgen allerede klokken halv syv, før jeg møder på arbejde}）。Katrine 在周日，Yuna 在周三晚。
  \item \textbf{7. A (Katrine)}：年费只需三百克朗，因为大家志愿轮流策划徒步活动（\textit{koster kun tre hundrede kroner om året, fordi alle skiftes til at planlægge turene frivilligt}）。David 反而说会费偏高（约两千克朗）。
\end{itemize}

\begin{transbox}
\small
\textbf{\color{NCDTeal}[Opgave 4 三人自述中文精译]}\\
\textbf{A. 卡特琳娜}：三年前刚搬到锡尔克堡不久，我就加入了当地的自然协会。那时我在城里几乎谁也不认识，不想一个人待在公寓里，很想和大家一起走进森林。每隔一个周日，我们十五到二十个人聚在一起徒步八到十公里。有时会有生物学家同行，沿途讲解植物和鸟类，之后我们还会架起篝火煮热汤或咖啡。会员年龄在二十到七十岁之间，因此能和各种各样的人交流。每年会费只需三百克朗，因为大家会自愿轮流策划路线。这对我来说是认识当地自然风光与周边地区最好的方式。\\
\textbf{B. 大卫}：一年半前我加入了峡湾边的赛艇俱乐部，因为整天坐在屏幕前让我腰背疼痛。一开始我必须参加四晚的必修入门课，学习水上技术与安全规则后才允许独自出海划艇。刚开始手臂真的挺吃力的！现在我每周二、周四早上六点半上班前就和固定的四人小组一起划艇。那个时段水面非常平静，能为我一整天注入充沛精力。会费略高，一年大约两千克朗，不过所有赛艇和救生衣都包含在内。夏天俱乐部还会组织去另一座岛的周末短途旅行。\\
\textbf{C. 优娜}：两年前我从日本来到奥胡斯时，很想在下班后做些有创意的事，同时在轻松的氛围里练习丹麦语。一位同事带我去了一家本地合唱团的排练夜，我一下子就迷上了。参加这里不需要会看乐谱也不用参加入学试音，所以人人都能加入。我们每周三晚上在文化馆排练，既唱丹麦歌曲也唱英文流行曲。最初几次排练时，我有点跟不上语速较快的丹麦歌词，但其他团员都耐心地帮我纠正发音。每年两次我们会在养老院或教堂举办小型演出。合唱团给了我强烈的归属感，如今好几位团员都成了我的挚友。
\end{transbox}

\paragraph*{五、核心高频词表（Gloser｜15 个必背 B1 屈折词条）}

\begin{glosertable}
名词 & en vagt & \textit{(vagten, vagter, vagterne)} & 值班，班次（如志愿者轮班） \\
名词 & et armbånd & \textit{(armbåndet, armbånd, armbåndene)} & 手环，入场腕带（单复数同形） \\
动词 & at overveje & \textit{(overvejer, overvejede, har overvejet)} & 考虑，斟酌 \\
副词 & alligevel & （句首触发 V2；句中置于中心状语位） & 仍然，尽管如此 \\
动词 & at fordele & \textit{(fordeler, fordelte, har fordelt)} & 分配，分摊 \\
名词 & en støj & \textit{(støjen, ingen flertal)} & 噪音，嘈杂声 \\
名词 & et lokale & \textit{(lokalet, lokaler, lokalerne)} & 店铺场地，房间，空间 \\
动词 & at opstå & \textit{(opstår, opstod, er opstået)} & 出现，产生（完成时用 \textit{er}！） \\
动词 & at fortryde & \textit{(fortryder, fortrød, har fortrudt)} & 后悔（不规则强变化！） \\
名词 & en opsparing & \textit{(opsparingen, opsparinger, -erne)} & 储蓄，积蓄 \\
名词 & en forening & \textit{(foreningen, foreninger, foreningerne)} & 协会，社团，俱乐部 \\
名词 & et kontingent & \textit{(kontingentet, kontingenter, -erne)} & 会费，会员费 \\
名词 & et fællesskab & \textit{(fællesskabet, fællesskaber, -ne)} & 共同体，社群归属感 \\
形容词 & pålidelig & \textit{(pålidelig, pålideligt, pålidelige)} & 可靠的，值得信赖的 \\
动词 & at anbefale & \textit{(anbefaler, anbefalede, har anbefalet)} & 推荐，建议 \\
\end{glosertable}

% =====================================================================
% 4.3 写作范文
% =====================================================================
\subsubsection*{4.3 写作双任务双档高分示范与点评（Skrivning Modeller）}

\begin{emailbox}{Opgave 1 范文一：B1 稳过版（109 词｜清晰覆盖 4 要点 · 句式自然）}
\textbf{Anbefaling af Byens Cykler}

Jeg vil gerne anbefale cykelhandleren \textit{Byens Cykler}, som ligger på Søndergade 42 i Aarhus lige ved biblioteket.

I sidste uge punkterede mit bagdæk på vej hjem fra arbejde, og mine bremser virkede heller ikke ordentligt. Jeg kørte derfor cyklen ind til værkstedet om morgenen uden at have bestilt tid.

Servicen var helt fantastisk. Mekanikeren tog imod mig med et smil og forklarede tydeligt, hvad der skulle laves, og hvad det ville koste. Allerede samme eftermiddag var cyklen færdig, og de havde endda smurt kæden gratis.

Jeg vil klart anbefale \textit{Byens Cykler} til alle i byen, der ønsker hurtig, ærlig og venlig hjælp til deres cykel.

\vspace{0.3em}
\rule{\linewidth}{0.4pt}\\[0.2em]
\small\textbf{\color{NCDTeal}[考官点评 · 109 ord]}：完整覆盖全部 4 项要求（地址、维修项目、服务优点、推荐对象）。语法亮点包括：\ding{172} 关系从句主语标配 \textit{hvad \textbf{der} skulle laves}；\ding{173} 句首状语倒装 \textit{Allerede samme eftermiddag \textbf{var} cyklen færdig}；\ding{174} 过去完成时 \textit{havde smurt} 补充背景；\ding{175} 不定式介词短语 \textit{uden at have bestilt}。
\end{emailbox}

\begin{emailbox}{Opgave 1 范文二：B1+ 高分示范版（141 词｜被动语态 + 介词接从句 + 高级衔接）}
\textbf{Den bedste cykelservice i byen}

Hvis du leder efter et pålideligt cykelværksted, kan jeg på det varmeste anbefale \textit{Byens Cykler}, som ligger centralt på Søndergade 42 i Odense, fem minutter fra banegården.

Forleden opdagede jeg, at gearet på min cykel var gået i stykker, og samtidig trængte både kæden og forlygten til at blive skiftet. \da{Selvom jeg ikke havde aftalt} nogen tid på forhånd, \dabold{blev cyklen taget ind} med det samme.

\da{Grunden til, at} jeg er så tilfreds med servicen, er den ærlige rådgivning. Inden reparationen gik i gang, fik jeg en klar pris, så jeg ikke \dabold{blev overrasket} ved kassen. Desuden sendte værkstedet en sms efter fire timer, og cyklen kørte som ny.

Jeg vil især anbefale stedet til pendlere og studerende, som er afhængige af deres cykel i hverdagen og sætter pris på god kvalitet til en fair pris.

\vspace{0.3em}
\rule{\linewidth}{0.4pt}\\[0.2em]
\small\textbf{\color{NCDNavy}[高分亮点 · 141 ord]}：\ding{172} 熟练运用 \textit{blive} 被动描述具体维修过程（\textit{blive skiftet / blev taget ind / blev overrasket}）；\ding{173} 让步从句内否定前置 + 主句倒装（\textit{Selvom jeg \textbf{ikke} havde aftalt..., \textbf{blev} cyklen...}）；\ding{174} 精准使用 B1 评价结构 \textit{Grunden til, at... er...} 与 \textit{sætter pris på}。
\end{emailbox}

\begin{emailbox}{Opgave 2 范文一：B1 稳过版（112 词｜严格 $\ge$ 90 词 · 回答全部 4 问）}
Kære Henrik

Tak for din søde mail!

Ja, det gik rigtig godt til jobsamtalen i torsdags, og jeg blev tilbudt jobbet allerede i går! Jeg er selvfølgelig utrolig glad.

Stillingen er som logistikmedarbejder på et mindre lager. Arbejdet handler om at pakke ordrer, tale med chauffører og holde styr på varerne i systemet.

Samtalen varede fyrre minutter og var meget afslappet. Jeg talte med chefen og en medarbejder. De var meget venlige og spurgte til min erfaring og mit dansk.

Før samtalen havde jeg forberedt mig grundigt på dansk. Jeg læste om firmaet, skrev små noter og øvede svarene højt med min nabo. Held og lykke med din jobsøgning!

Kærlig hilsen\\
Klara

\vspace{0.3em}
\rule{\linewidth}{0.4pt}\\[0.2em]
\small\textbf{\color{NCDTeal}[考官点评 · 112 ord]}：远超 90 词底线，称呼落款得体。顺畅使用过去被动 \textit{blev tilbudt} 与过去完成时 \textit{havde forberedt mig} 区分“面试前准备”与“面试当天经过”，完整回答 Henrik 的全部四个问题，信息密度充实。
\end{emailbox}

\begin{emailbox}{Opgave 2 范文二：B1+ 高分示范版（138 词｜介词 + om 从句 + 对比连接）}
Kære Henrik

Tak for din besked -- det er dejligt at høre, at du også søger arbejde!

Det gik over al forventning til min jobsamtale i sidste uge, for i går ringede de og tilbød mig stillingen.

Jobbet er som kontorassistent i en grøn virksomhed, hvor mine vigtigste opgaver bliver at svare på kundehenvendelser, planlægge møder og hjælpe med mindre regnskaber.

\da{Selvom jeg i starten var lidt spændt på, om} mit danske sprog var godt nok, \dabold{blev samtalen} hurtigt tryg og uformel. Vi drak kaffe, og de to ledere lagde vægt på min motivation frem for perfekte endelser.

For at blive klar havde jeg brugt weekenden på at forberede mig. \da{Ud over at undersøge} hjemmesiden trænede jeg spørgsmål med en dansk veninde, så jeg kendte de vigtigste fagord. Skriv endelig, hvis du vil låne mine noter!

Mange hilsner\\
Klara

\vspace{0.3em}
\rule{\linewidth}{0.4pt}\\[0.2em]
\small\textbf{\color{NCDNavy}[高分亮点 · 138 ord]}：\ding{172} 精准示范介词接间接疑问从句（\textit{spændt på, om mit danske sprog var godt nok}）；\ding{173} 高级短语 \textit{over al forventning}、\textit{lagde vægt på... frem for...} 与 \textit{Ud over at...} 极大提升书面表达成熟度，且字数控制在考场 30 分钟黄金区间。
\end{emailbox}

% =====================================================================
% 4.4 口语示范稿与四图双人对话
% =====================================================================
\subsubsection*{4.4 口语示范逐字稿、核心 IPA 与四图双人互动全真对话}

\begin{transbox}
\small
\textbf{\color{NCDNavy}[Mindmap 3A 口语陈述范文：Min skole / uddannelse｜152 词 · 约 1.5 分钟]}\\
\textit{Jeg vil gerne fortælle lidt om min uddannelse \ipa{ˈuðˌdanˀəlsə}. Jeg studerer multimediedesign på et erhvervsakademi \ipa{ɛɐˈvɛːˀɐsˌɑkʰadəˌmiˀ} i København, hvor jeg går på andet år. Det er en offentlig uddannelse, og der går cirka to hundrede studerende \ipa{sd̥uˈdeːˀɐnə} på min årgang. Jeg møder fra mandag til fredag mellem klokken ni og femten. Hver morgen cykler jeg langs kanalen i tyve minutter, men hvis det regner kraftigt, tager jeg metroen i stedet for. I undervisningen \ipa{ˈʌnɐˌviːˀsneŋ} arbejder vi både med grafisk design, kodning og store gruppeprojekter \ipa{ˈɡʁubəˌpʰʁoˌjɛɡdɐ}. Jeg kan bedst lide de praktiske opgaver, mens det sværeste for mig stadig er at skrive lange rapporter på dansk. Vores lærere er meget uformelle og hjælpsomme, og mit hold består af både danskere og internationale studerende. Jeg mødes fast med en lille studiegruppe \ipa{ˈsd̥uˀdjəˌɡʁubə}, hvor vi støtter hinanden. Alt i alt er jeg rigtig glad for uddannelsen, fordi den er jobrelevant, selvom det bliver ret travlt før eksamen \ipa{ɛɡˈsæːmən}.}
\end{transbox}

\paragraph*{Mindmap 3A Del 2 考官追问实战示范对话（4 题全解）：}
\dialogturn{1}{Eksaminator}{Du siger, at I laver mange gruppeprojekter --- kan du bedst lide at arbejde alene eller i en gruppe? Hvorfor?}{gruppeprojekter \ipa{ˈɡʁubəˌpʰʁoˌjɛɡdɐ} · alene \ipa{aˈleːnə}}{你说你们做很多小组项目——你更喜欢独自工作还是小组合作？为什么？}
\dialogturn{2}{Kursist}{Jeg kan faktisk bedst lide at arbejde i en gruppe, fordi man får flere gode idéer sammen. Når man er usikker på et dansk fagord, kan ens studiekammerater også hjælpe med det samme.}{idéer \ipa{iˈdeːˀɐ} · studiekammerater \ipa{ˈsd̥uˀdjəˌkʰamɑˌʁɑːˀdɐ}}{其实我更喜欢小组合作，因为大家一起能碰撞出更多好点子。当对某个丹麦语专业词拿不准时，同学也能立刻帮忙。}
\dialogturn{3}{Eksaminator}{Du nævner, at skriftlige rapporter på dansk er det sværeste --- hvad gør du for at blive bedre til det?}{skriftlige \ipa{ˈsɡʁefdliə} · rapporter \ipa{ʁaˈpʰʌdɐ}}{你提到用丹麦语写书面报告是最难的——你做些什么来提升呢？}
\dialogturn{4}{Kursist}{Jeg læser mange danske artikler om design, og når vi skriver i gruppen, beder jeg ofte mine danske kammerater om at rette mine sætninger. På den måde lærer jeg hurtigt af mine fejl.}{artikler \ipa{ɑˈtˢiɡlɐ} · fejl \ipa{ˈfɑi̯ˀl}}{我读很多关于设计的丹麦语文章，而且小组写作时常请丹麦同学帮我修改句子。这样我能很快从错误中吸取经验。}
\dialogturn{5}{Eksaminator}{Hvordan er forholdet mellem lærere og studerende i Danmark sammenlignet med dit hjemland?}{forholdet \ipa{ˈfɒːˌhʌlˀəð} · sammenlignet \ipa{ˈsamənˌliˀnəð}}{与你的祖国相比，丹麦师生之间的关系怎么样？}
\dialogturn{6}{Kursist}{Her er tonen meget mere uformel. Vi kalder lærerne ved fornavn og drikker kaffe sammen i pauserne. Det gør, at man tør stille dumme spørgsmål uden at være nervøs.}{uformel \ipa{ˈufʌˌmɛlˀ} · fornavn \ipa{ˈfɒːˌnɑu̯ˀn}}{这里的氛围要非正式得多。我们直呼老师名字，课间还一起喝咖啡，这让人敢于大胆提问而不觉得紧张。}
\dialogturn{7}{Eksaminator}{Hvilken slags job drømmer du om at få, når du er færdig med uddannelsen?}{drømmer \ipa{ˈdʁœmɐ} · færdig \ipa{ˈfɛːˀdi}}{毕业后你梦想找一份什么样的工作？}
\dialogturn{8}{Kursist}{Jeg håber at få et job i et mellemstort designbureau, hvor jeg både kan lave hjemmesider og brugervenlige apps --- og helst på en arbejdsplads, hvor vi taler dansk til hverdag.}{designbureau \ipa{diˈsɑi̯nbyˌʁo} · brugervenlige \ipa{ˈbʁuːɐˌvɛnliə}}{我希望在一家中型设计事务所工作，既做网站也做用户友好的应用——最好是在日常讲丹麦语的职场。}

\begin{transbox}
\small
\textbf{\color{NCDNavy}[Mindmap 3B 口语陈述范文：Mit frivillige arbejde｜157 词 · 约 1.5 分钟]}\\
\textit{Nu vil jeg gerne fortælle om mit frivillige arbejde \ipa{ˈfʁiˌviliə ˈɑːbɑi̯ˀdə}. Hver anden torsdag aften hjælper jeg til i et lokalt folkekøkken, som hedder Naboens Bord, og det ligger i et kulturhus i midtbyen. Jeg er der fra klokken halv fem til halv ni, så jeg bruger cirka otte timer om måneden på det. Mine arbejdsopgaver \ipa{ˈɑːbɑi̯ˀdsˌʌbˌɡæːvɐ} er ret varierede: Først skærer jeg grøntsager og hjælper kokken, og bagefter dækker jeg langborde og tager imod gæsterne ved døren. Til sidst rydder vi op i køkkenet sammen. Vi er et fast hold på ti frivillige i alle aldre --- både naboer og unge studerende --- og der er altid en rigtig god holdånd \ipa{ˈhʌlˌʌnˀ}. Det har givet mig et stort udbytte \ipa{ˈuðˌbydə}, for jeg øver hverdagssnak på dansk, har fået nye venner og lært meget om dansk madkultur. Selvom man bliver lidt træt i benene af at stå op længe, synes jeg, at arbejdet er utrolig meningsfuldt \ipa{ˈmeːneŋsˌfult} og giver mig ny energi.}
\end{transbox}

\paragraph*{Mindmap 3B Del 2 考官追问实战示范对话（4 题全解）：}
\dialogturn{1}{Eksaminator}{Hvordan fandt du frem til dette folkekøkken, og hvorfor valgte du at blive frivillig?}{folkekøkken \ipa{ˈfʌlɡəˌkʰøɡən} · valgte \ipa{ˈvalˀdə}}{你是怎么找到这家社区食堂的？为什么选择当志愿者？}
\dialogturn{2}{Kursist}{En kollega fortalte mig om stedet, og jeg havde lyst til at gøre noget socialt uden for mit arbejde. Mad samler jo folk, så det føltes som en helt naturlig måde at møde lokale danskere på.}{socialt \ipa{soˈɕæːˀld} · samler \ipa{ˈsamlɐ}}{一位同事告诉我的，我也想在工作之余做点公益社交。美食总能把大家聚在一起，所以感觉这是认识本地丹麦人非常自然的方式。}
\dialogturn{3}{Eksaminator}{Hvad kan du bedst lide af dine opgaver i køkkenet eller ved bordene?}{opgaver \ipa{ˈʌbˌɡæːvɐ} · bordene \ipa{ˈboːˀɐnə}}{在厨房或餐桌旁的各项任务中，你最喜欢哪一项？}
\dialogturn{4}{Kursist}{Jeg elsker mest at tage imod gæsterne og servere maden, fordi man så får mange små samtaler i gang. Det er lidt mindre sjovt at vaske de store gryder op til sidst, men når vi gør det sammen med musik, går det hurtigt.}{servere \ipa{sɛɐˈveːˀɐ} · gryder \ipa{ˈɡʁyːðɐ}}{我最喜欢迎接客人和上菜，因为那样能和不少人聊起来。最后刷大锅没那么有趣，不过大家放着音乐一起干就很快。}
\dialogturn{5}{Eksaminator}{Hvad har det frivillige arbejde betydet for dit danske sprog og dit netværk?}{betydet \ipa{beˈtˢyːˀðəð} · netværk \ipa{ˈnɛdˌvɛɐɡ}}{这段志愿经历对你的丹麦语和人脉圈有什么意义？}
\dialogturn{6}{Kursist}{Det har betydet rigtig meget! I begyndelsen var jeg genert, men nu tør jeg tale dansk om alt muligt ved bordene. Og to af de andre frivillige inviterede mig faktisk til fødselsdag i sidste måned.}{genert \ipa{ɕeˈneːˀɐd} · inviterede \ipa{enˈviˌtˢeːˀɐðə}}{意义非常重大！起初我很害羞，但现在在餐桌旁什么话题都敢用丹麦语聊。上个月还有两位志愿者朋友邀请我参加了生日聚会。}
\dialogturn{7}{Eksaminator}{Synes du, at det er let eller svært at passe frivilligt arbejde ved siden af hverdagen?}{passe \ipa{ˈpʰasə} · hverdagen \ipa{ˈvɛːˀˌdæːˀən}}{你觉得在日常工作生活之余兼顾志愿服务容易还是困难？}
\dialogturn{8}{Kursist}{Det er faktisk ret nemt, når det kun er hver anden torsdag. Fordi vi planlægger vagterne en måned i forvejen, kan man sagtens få det til at passe med familie og arbejde.}{vagterne \ipa{ˈvaɡdɐnə} · i forvejen \ipa{i ˈfɒːˌvɑi̯ˀən}}{隔周周四去一次其实挺容易的。因为我们提前一个月排好班次，完全能和家庭及工作协调好。}

\paragraph*{Mundtlig Opgave 2（EMNE: Indkøb og forbrug）Del 1 双人四图全真互动示范：}

\dialogturn{1}{Kursist A}{Skal vi kigge på det første billede med supermarkedet? Jeg handler selv i supermarked tre gange om ugen, fordi det ligger lige om hjørnet og har alle dagligvarer på ét sted. Hvad synes du om at købe ind i supermarkedet?}{supermarked \ipa{ˈsuːˀpʰɐˌmɑːɡəð} · dagligvarer \ipa{ˈdɑuliˌvɑːɐ}}{我们看看第一张超市图吧？我自己每周去超市买三次，因为它就在拐角而且一站买齐日用品。你觉得去超市买菜怎么样？}
\dialogturn{2}{Kursist B}{Jeg kan egentlig også godt lide det til hverdag, især når man vil vælge frisk frugt og kød selv. Men jeg synes tit, der er lange køer klokken fem om eftermiddagen. Foretrækker du i stedet at handle på nettet, ligesom på billede to?}{foretrækker \ipa{ˈfɒːˌtˢʁɛɡɐ} · nettet \ipa{ˈnɛdəð}}{平时我也挺喜欢的，尤其是想自己挑新鲜水果和肉的时候。不过下午五点经常排长队。你会更喜欢像图二那样网购吗？}
\dialogturn{3}{Kursist A}{Ja, nethandel bruger jeg rigtig meget til tunge ting som havregryn og vaskepulver, og også til tøj, fordi det sparer en masse tid. Til gengæld kan det være irriterende at sende tøj tilbage, hvis størrelsen ikke passer. Handler du ofte online?}{nethandel \ipa{ˈnɛdˌhanˀəl} · størrelsen \ipa{ˈsd̥œːɐlsən}}{是的，买燕麦片、洗衣粉等重物或衣服时我常用网购，因为非常省时间。但衣服尺码不合适还得退货挺烦的。你经常网购吗？}
\dialogturn{4}{Kursist B}{Kun nogle gange til elektronik eller bøger, for jeg bliver faktisk lidt træt af alle de papkasser. Men jeg elsker billede tre med loppemarked og genbrugsbutik! Jeg køber næsten alle mine tallerkener og lamper brugt, fordi det både er billigt og godt for klimaet. Hvad mener du om genbrug?}{loppemarked \ipa{ˈlʌbəˌmɑːɡəð} · klimaet \ipa{ˈkʰliːˌmæːˀəð}}{我只偶尔网购电子产品或书，因为收一堆纸箱挺烦的。但我超爱图三的跳蚤市场和二手店！我几乎所有盘子和台灯都买二手的，既便宜又环保。你怎么看二手购物？}
\dialogturn{5}{Kursist A}{Jeg er enig i, at genbrug er super til møbler og bøger, selvom det tager lidt længere tid at finde den rigtige ting. Hvad så med det sidste billede fra en gårdbutik eller et marked? Har du prøvet det?}{enig \ipa{ˈeːni} · gårdbutik \ipa{ˈɡɔːbuˌtˢiɡ}}{我同意买二手家具和书很棒，尽管挑到合适物件要多花点时间。那最后一张农场直营店或集市图呢？你去过吗？}
\dialogturn{6}{Kursist B}{Ja, om sommeren cykler vi tit forbi en lille gårdbutik på landet. Grøntsagerne smager meget bedre end i supermarkedet, og man støtter lokale landmænd. Dog er priserne lidt højere, så det er mest til weekenden. Er du også villig til at betale ekstra for friske varer?}{landmænd \ipa{ˈlanˌmɛnˀ} · villig \ipa{ˈvili}}{去过，夏天我们常骑车路过乡下小农场商店。蔬菜比超市好吃得多，还能支持本地农户。不过价格略贵，所以主要周末去。你也愿意为新鲜食材多花钱吗？}
\dialogturn{7}{Kursist A}{Ja, til gæster og i weekenden gør jeg gerne det, fordi kvaliteten bare er i top. Men til daglig er supermarkedet og et godt tilbud på nettet stadig det mest praktiske for mig.}{kvaliteten \ipa{kʰvaliˈtˢeːˀdən} · praktiske \ipa{ˈpʰʁɑɡtˢisɡə}}{是的，待客或周末我愿意，因为品质真是一流。但工作日还是超市和划算的网购最方便我。}

\paragraph*{Mundtlig Opgave 2 Del 2（考官个人深入追问 3 题高分示范）：}
\dialogturn{1}{Eksaminator (Q1)}{Hvordan kan du bedst lide at købe ind --- og hvorfor?}{hvordan \ipa{vʌˈdan} · bedst \ipa{ˈbɛsd}}{你最喜欢以哪种方式购物——为什么？}
\dialogturn{2}{Kursist}{Jeg kan bedst lide en blanding af supermarked i hverdagene og gårdbutik i weekenden. I supermarkedet går det hurtigt efter arbejde, mens jeg på et lokalt marked virkelig nyder den gode kvalitet og atmosfæren.}{blanding \ipa{ˈblaneŋ} · atmosfæren \ipa{admosˈfɛːˀɐn}}{我最喜欢工作日逛超市、周末去农场店结合。下班后在超市采购很快捷，而周末在本地集市能好好享受高品质食材和热闹氛围。}
\dialogturn{3}{Eksaminator (Q2)}{Hvad er vigtigst for dig, når du handler i hverdagen --- pris, kvalitet, tid, klima eller økologi?}{økologi \ipa{økʰoloˈɡiˀ} · vigtigst \ipa{ˈveɡtˢisd}}{平时买东西时对你最重要的是什么——价格、质量、时间、气候环保还是有机？}
\dialogturn{4}{Kursist}{Tid og kvalitet betyder mest for mig, men jeg tænker også meget på økologi og madspild. Derfor planlægger jeg ofte en madplan for hele ugen og køber danske, økologiske grøntsager, når de er i sæson.}{madspild \ipa{ˈmaðˌsbilˀ} · sæson \ipa{sɛˈsʌŋ}}{时间和质量对我最重要，但我也非常重视有机和减少食物浪费。因此我常制定一周食谱，专门买当季的丹麦有机蔬菜。}
\dialogturn{5}{Eksaminator (Q3)}{Kan du fortælle lidt om et godt eller dårligt køb, du har lavet for nylig?}{for nylig \ipa{fʌ ˈnyːli} · køb \ipa{ˈkʰøːˀb}}{能讲讲你最近一次很满意或很糟糕的购物经历吗？}
\dialogturn{6}{Kursist}{For to uger siden gjorde jeg et rigtig godt køb på et loppemarked i Aarhus. Jeg fandt et næsten nyt spisebord i massivt egetræ til kun fem hundrede kroner. Sælgeren hjalp endda med at bære det ned til min vens bil, så jeg sparede flere tusinde kroner.}{massivt \ipa{maˈsiu̯ˀd} · egetræ \ipa{ˈɑi̯əˌtˢʁæˀ}}{两周前我在奥胡斯跳蚤市场捡了个大漏。我只花五百克朗买到一张几乎全新的实木橡木餐桌。卖家还帮忙抬到我朋友车上，让我省下了好几千克朗。}
